% Copyright page
\thispagestyle{empty}
% \null\vfill

\begin{center}
\includegraphics[width=0.2\linewidth,keepaspectratio]{YHK_Color_OutOfTheBox_tight} \\[1.5em]

\textbf{\Huge मेन्टल मॉडेल्स}\\ [0.5em]
{\small(विविध क्षेत्रातील `मेंटल मॉडेल्स'ची तोंडओळख )}\\[0.5em]

लेखक: \textbf{{\large डॉ. योगेश हरिभाऊ कुलकर्णी}}\\[1.5em]
\end{center}

\vspace{3.5em}

\begin{flushleft}

प्रकाशक: \ldots \\ %डॉ. योगेश हरिभाऊ कुलकर्णी (self-published)
पत्ता:  \ldots \\ %पाषाण ,  पुणे ८
फोन:  \ldots \\ % +91 9890251406
ईमेल: \ldots \\[1.5em] %yogeshkulkarni@yahoo.com

\vspace{3.5em}

प्रथम आवृत्ती: \ldots\\[0.5em] %  २०२६ 

% ISBN: [ISBN नंबर]\\[0.5em]

कॉपीराइट © \dots \\[0.5em] % २०२५ डॉ. योगेश हरिभाऊ कुलकर्णी

{\textit{सर्व हक्क राखीव. या पुस्तकाचा कोणताही भाग प्रकाशकाच्या लेखी परवानगीशिवाय कोणत्याही स्वरूपात पुनर्मुद्रित किंवा पुनर्प्रकाशित करता येणार नाही.  या पुस्तकात व्यक्त केलेली मते लेखकाची व्यक्तिगत आहेत.}}\\[1.5em]

{\large Legal Notice:}\\
{\textit{All rights reserved. No part of this publication may be reproduced, distributed, or transmitted in any form or by any means, including photocopying, recording, or other electronic or mechanical methods, without the prior written permission of the publisher, except in the case of brief quotations embodied in critical reviews and certain other noncommercial uses permitted by copyright law.}}
\end{flushleft}
\vfill\null
\clearpage

\begin{dedication}
`मेंटल मॉडेल्स 'यांना प्रसिद्धी देणाऱ्या चार्ली मंगर आणि शॉन पेरिश यांना समर्पित  
\end{dedication}

\clearpage

\chapter*{पुस्तकाविषयी }
जगाकडे पाहण्याची प्रत्येकाची पद्धत वेगळी असते. एखाद्या समस्येचे आकलन, निर्णय घेण्याची शैली किंवा परिस्थितीचा अर्थ लावण्यामागे काही विशिष्ट विचारचौकटी कार्यरत असतात. आधुनिक भाषेत या विचारचौकटींना “मेंटल मॉडेल्स” असे म्हटले जाते.
हे पुस्तक अशाच विविध मेन्टल मॉडेल्सची ओळख करून देणारे आहे. व्यवस्थापन, मानसशास्त्र, अर्थशास्त्र, विज्ञान आणि मानवी वर्तन यांमधील मूलभूत कल्पना साध्या उदाहरणांतून मांडत हे पुस्तक वाचकाला अधिक सजग आणि व्यापक विचार करण्यास प्रवृत्त करते. विद्यार्थी, व्यावसायिक, उद्योजक किंवा जिज्ञासू वाचक,  प्रत्येकासाठी हे पुस्तक उपयुक्त ठरू शकते.
या पुस्तकातील लेख यापूर्वी दै. 'सकाळ 'मधील ‘तिसरा मेंदू’ या सदरातून प्रकाशित झालेले असून, आवश्यक त्या संपादनासह येथे एकत्रित स्वरूपात सादर करण्यात आले आहेत. काही लेख 'मीडियम' या संकेतस्थळावरील ‘देसी स्टॅक’ या प्रकाशनातही उपलब्ध आहेत. वाचकांना एकसंध आणि संदर्भासाठी उपयुक्त असा संग्रह मिळावा, हा या पुस्तकामागील प्रमुख हेतू आहे.
या आवृत्तीत पुस्तकातील चित्रे नव्याने साकारण्यात आली आहेत. लेखकाने काढलेल्या मूळ संकल्पनाचित्रांना कृत्रिम बुद्धिमत्ता (AI) तंत्रज्ञानाच्या साहाय्याने दृश्यरूप देण्यात आले आहे. तसेच प्रत्येक प्रकरणाच्या शेवटी दिलेला इंग्रजीतील संक्षिप्त सारांश वाचकांना या विषयावरील जागतिक साहित्याकडे जाण्यासाठी ‘सेतू’ म्हणून उपयोगी ठरेल.
हे पुस्तक तयार उत्तरे देत नाही; उलट अधिक चांगले प्रश्न विचारण्याची दृष्टी विकसित करण्याचा प्रयत्न करते. विचारांच्या नव्या चौकटींची ही ओळख वाचकांच्या निर्णयप्रक्रियेत अधिक स्पष्टता, सजगता आणि व्यापकता आणेल, हीच अपेक्षा.

% \chapter*{शीर्षकाविषयी }
% `मन:प्रारूपे' या पुस्तकाचे शीर्षक स्वतःतच एक लहानसं विधान आहे. विचारांच्या जगाचे दार उघडणारे, आणि मानवी बुद्धीच्या प्रवासाला दिशादर्शक ठरणारे. येथे ‘मन’ हा शब्द केवळ भावना किंवा विचारांचे स्थान म्हणून नाही, तर मानवी जाणिवा, समज, निर्णयक्षमता आणि आकलन यांचे केंद्र म्हणून वापरला आहे. मन म्हणजे आपण जग अनुभवतो त्या सर्व प्रक्रियांचा पाया, आणि त्यातूनच उगम पावतात आपले निर्णय व निष्कर्ष.

% ‘प्रारूपे’ हा शब्द या संदर्भात विशेष अर्थपूर्ण आहे. प्रारूप म्हणजे एखाद्या वस्तूचे, घटनेचे किंवा विचाराचे आंतरिक रूप, त्याची चौकट, त्याचे कार्य समजून घेण्याची रचना. इंग्रजीतील ‘मॉडेल्स’ किंवा ‘फ्रेमवर्क्स’ यांची मराठी प्रतिमा म्हणून ‘प्रारूपे’ हा शब्द; तो विचारांच्या मुळावर बोट ठेवतो. ही प्रारूपे म्हणजे आपल्या मनात नकळतपणे तयार होणारे नमुने, जग पाहण्याचे, समस्या विश्लेषित करण्याचे आणि निर्णय घेण्याचे मार्गदर्शक नकाशे.

% या शीर्षकाने पुस्तकाचा आशय थेट सुचवला आहे. ते वाचकाला सांगते की पुढील पानांमध्ये आपण त्या मानसिक नकाशांचा अभ्यास करणार आहोत, जे आपल्या वर्तनाला दिशा देतात, पूर्वग्रहांना जन्म देतात, आणि समजुतींची मर्यादा किंवा व्यापकता ठरवतात.  

\chapter*{मनोगत}
आपण जगाकडे ज्या नजरेने पाहतो, तेच आपले जग बनत जाते. प्रत्येक व्यक्ती विचार करते, निर्णय घेते, मत बनवते; पण त्या सर्व प्रक्रियांमागे काही अदृश्य विचारचौकटी कार्यरत असतात. कधी त्या अनुभवांतून तयार झालेल्या असतात, कधी शिक्षणातून, तर कधी नकळत स्वीकारलेल्या समजुतींतून. या विचारांच्या चौकटी अधिक स्पष्ट, व्यापक आणि सजगपणे समजून घेण्याचा प्रयत्न म्हणजे हे पुस्तक.

“मेंटल मॉडेल्स” हा विषय मला आकर्षित करण्यामागचे कारण म्हणजे त्याची व्यावहारिकता. अनेक तत्त्वज्ञानाशी संबंधित किंवा बौद्धिक संकल्पना वाचताना प्रभावी वाटतात; परंतु दैनंदिन जीवनात त्यांचा उपयोग कसा करायचा, हा प्रश्न कायम राहतो. मेन्टल मॉडेल्स मात्र वेगळी आहेत. ती केवळ विचार करण्याबद्दल बोलत नाहीत, तर अधिक चांगल्या प्रकारे विचार कसा करावा यासाठी साधने देतात. निर्णय घेणे, चुका समजून घेणे, पर्यायांचा विचार करणे, दीर्घकालीन परिणाम ओळखणे किंवा स्वतःच्या विचारांमधील मर्यादा शोधणे, या सर्व गोष्टींसाठी ही मॉडेल्स उपयुक्त ठरतात.

हे लेखन करताना माझा प्रयत्न एकच होता की अवघड वाटणाऱ्या संकल्पना शक्य तितक्या साध्या भाषेत मांडणे. मराठीत तंत्रज्ञान, विज्ञान किंवा विचारप्रणाली यांसारख्या विषयांवर लेखन करताना अनेकदा भाषा अतिशय तांत्रिक बनते. मला मात्र हे पुस्तक संवादासारखे वाटावे, अशी इच्छा होती. वाचकांनी एखाद्या मित्राशी गप्पा माराव्यात तसे हे विचार वाचावेत, त्यावर थांबून विचार करावा आणि आपल्या जीवनाशी त्यांची सांगड घालावी, एवढीच अपेक्षा आहे.

या पुस्तकातील लेख दै. सकाळ मधील ‘तिसरा मेंदू’ या सदरातून प्रथम वाचकांसमोर आले. त्या लेखांना मिळालेला प्रतिसाद आणि वाचकांकडून आलेल्या प्रतिक्रिया यांनी त्या लेखांचे पुस्तक स्वरूपात संकलन करण्याची प्रेरणा दिली. या प्रवासात सातत्याने प्रोत्साहन आणि सहकार्य दिल्याबद्दल 'सकाळ'चे निरंजन आगाशे यांचे मनःपूर्वक आभार.

हे पुस्तक प्रकाशित करण्यासाठी पुढाकार घेतल्याबद्दल 'सकाळ प्रकाशन'चे आशुतोष रामगीर, दीपाली चौधरी, ऐश्वर्या कुमठेकर आणि त्यांच्या संपूर्ण टीमचा मी विशेष ऋणी आहे. नवीन विषय मराठी वाचकांपर्यंत पोहोचवण्याची त्यांची तयारी आणि विश्वास यामुळेच हे पुस्तक साकार होऊ शकले.

माझ्या जडणघडणीत महत्त्वाची भूमिका असलेल्या माझे वडील श्री. हरिभाऊ कुलकर्णी आणि आई सौ. सुनंदा कुलकर्णी यांच्याप्रती मी मनःपूर्वक कृतज्ञता व्यक्त करतो. माझी पत्नी सौ. अंजली हिने प्रत्येक टप्प्यावर दिलेले प्रोत्साहन आणि खंबीर साथ माझ्यासाठी सतत प्रेरणादायी ठरली आहे. तसेच माझ्या मुली, अनुष्का आणि अनुजा, यांनी आयुष्यात निर्माण केलेला उत्साह, आनंद आणि सकारात्मक ऊर्जा यांमुळे अनेक नव्या कल्पनांना प्रेरणा मिळाली आणि हा लेखनप्रवास अधिक समृद्ध झाला.

हे पुस्तक वाचताना प्रत्येक वाचकाला किमान एक तरी असा विचार सापडावा, जो त्याच्या जगाकडे पाहण्याच्या दृष्टिकोनात छोटासा का होईना, बदल घडवेल हीच या लेखनप्रवासामागची खरी प्रेरणा आहे.

